\clearpage\pagestyle{otMain}\fontsize{11}{15}\selectfont\corrigesChapitre{11}{Probabilités, évènements et simulations}
\begin{blocLivre}{cor-11-1}\begin{corrige}{1}{Un dé}\demarche{Énumérer les issues favorables.}
Les issues favorables sont 2, 4 et 6, parmi six issues équiprobables. La probabilité vaut $3/6=\res{1/2}$.
\end{corrige}\end{blocLivre}
\begin{blocLivre}{cor-11-2}\begin{corrige}{2}{Des boules colorées}\demarche{Compter les boules plutôt que les couleurs.}
Le total vaut 10. $P(R)=5/10=\res{1/2}$ et $P(\overline R)=1-1/2=\res{1/2}$. Trois couleurs ne signifient pas trois issues équiprobables.
\end{corrige}\end{blocLivre}
\begin{blocLivre}{cor-11-3}\begin{corrige}{3}{Écrire des évènements}\demarche{Lister les issues de chaque évènement avant de les combiner.}
$A=\{2;4;6\}$, $B=\{3;6\}$. $\res{A\cap B=\{6\}}$, $\res{A\cup B=\{2;3;4;6\}}$ et $\res{\overline A=\{1;3;5\}}$. L’issue 6 n’est comptée qu’une fois dans la réunion.
\end{corrige}\end{blocLivre}
\begin{blocLivre}{cor-11-4}\begin{corrige}{4}{Impossible ou certain}\demarche{Comparer l’évènement à toutes les issues.}
Obtenir 9 est impossible : probabilité $\res0$ ; cet évènement est l’ensemble vide $\varnothing$. Toutes les faces sont inférieures à 7 : l’évènement est certain, de probabilité $\res1$.
\end{corrige}\end{blocLivre}
\begin{blocLivre}{cor-11-5}\begin{corrige}{5}{Une fréquence expérimentale}\demarche{Ne pas confondre résultat observé et modèle.}
La fréquence est $46/100=\res{0{,}46}$. Cet écart à 0,5 peut se produire avec une pièce équilibrée ; ce seul résultat n’est pas une preuve. La fréquence dépend de l’expérience réalisée.
\end{corrige}\end{blocLivre}
\begin{blocLivre}{cor-11-6}\begin{corrige}{6}{La relation de la réunion}\demarche{Isoler $P(A\cup B)$ dans la relation.}
$P(A)=3/6=1/2$, $P(B)=2/6=1/3$ et $P(A\cap B)=1/6$. La relation donne $P(A\cup B)=1/2+1/3-1/6=\res{2/3}$. Contrôle : $A\cup B$ contient 4 issues sur 6, et $4/6=2/3$.
\end{corrige}\end{blocLivre}
\begin{blocLivre}{cor-11-7}\begin{corrige}{7}{Deux pièces}\demarche{Suivre toutes les branches de l’arbre.}
Les quatre issues équiprobables sont PP, PF, FP, FF. Exactement une face : PF et FP, d’où $\res{1/2}$. Au moins une face : PF, FP, FF, d’où $\res{3/4}$, ou $1-1/4$ par le contraire PP.
\label{fig-cor-11-7}\figureGeo[corrige]{proba-pieces}
\end{corrige}\end{blocLivre}
\begin{blocLivre}{cor-11-8}\begin{corrige}{8}{La somme de deux dés}\demarche{Compter des couples et non des sommes.}
Les 36 couples ordonnés sont équiprobables. Les six couples $(1;6),(2;5),(3;4),(4;3),(5;2),(6;1)$ conviennent : $\res{1/6}$. Les sommes elles-mêmes ne sont pas équiprobables.
\label{fig-cor-11-8}\figureGeo[corrige]{proba-des}
\end{corrige}\end{blocLivre}
\begin{blocLivre}{cor-11-9}\begin{corrige}{9}{Un « ou » exclusif}\demarche{Retirer de la réunion les issues communes.}
$A\cup B=\{2;4;5;6\}$ et $A\cap B=\{4;6\}$. Les issues cherchées sont 2 et 5 : probabilité $2/6=\res{1/3}$. Par le calcul : $P(A\cup B)-P(A\cap B)=4/6-2/6$.
\end{corrige}\end{blocLivre}
\begin{blocLivre}{cor-11-10}\begin{corrige}{10}{Au moins un six}\demarche{Utiliser la relation de la réunion pour éviter un double comptage.}
« Au moins un six » est $A\cup B$. $P(A)=P(B)=6/36$ et $A\cap B$ ne contient que $(6;6)$. Donc $P(A\cup B)=6/36+6/36-1/36=\res{11/36}$. Additionner $1/6+1/6$ aurait compté deux fois le couple $(6;6)$.
\label{fig-cor-11-10}\figureGeo[corrige]{proba-des}
\end{corrige}\end{blocLivre}
\begin{blocLivre}{cor-11-11}\begin{corrige}{11}{Deux tirages avec remise}\demarche{Numéroter les boules pour obtenir des issues équiprobables.}
Le tableau compte $5\times5=25$ couples équiprobables. Deux rouges : $3\times3=9$ couples, d’où $\res{9/25}$. Deux bleues : 4 couples. Même couleur : $9+4=13$ couples, d’où $\res{13/25}$, car les deux évènements sont incompatibles.
\label{fig-cor-11-11}\figureGeo[corrige]{proba-tableau-boules}
\end{corrige}\end{blocLivre}
\begin{blocLivre}{cor-11-12}\begin{corrige}{12}{Un jeu à points}\demarche{Distinguer gain brut, mise et bilan net.}
Le gain brut moyen vaut $8/6+4/6=\res2$ points. Il égale la mise : le bilan net moyen théorique est 0. Cela ne garantit ni un remboursement ni un gain à chaque partie ; un résultat individuel peut être $-2$, 2 ou 6 points nets.
\end{corrige}\end{blocLivre}
\begin{blocLivre}{cor-11-13}\begin{corrige}{13}{Un effectif inconnu}\demarche{Traduire le total avant le produit en croix.}
On écrit $6/(6+b)=2/5$. Les produits en croix donnent $30=12+2b$, donc $\res{b=9}$. Avec 15 boules, $6/15=2/5$, ce qui confirme le résultat.
\end{corrige}\end{blocLivre}
\begin{blocLivre}{cor-11-14}\begin{corrige}{14}{Un code à deux chiffres}\demarche{Préciser si la répétition est autorisée, puis utiliser la relation de la réunion.}
Il y a $4\times4=16$ codes équiprobables, dont $4\times3=12$ à chiffres distincts : $\res{3/4}$. Ensuite $P(\text{pair})=8/16$, $P(\text{début }0)=4/16$ et l’intersection contient 00 et 02 : $P=8/16+4/16-2/16=\res{5/8}$.
\end{corrige}\end{blocLivre}
\begin{blocLivre}{cor-11-15}\begin{corrige}{15}{Estimer sans promettre}\demarche{Comparer les écarts absolus sans en faire une règle monotone.}
Les fréquences sont 0,327 et 0,3354. Leurs écarts à $1/3$ valent environ 0,00633 et 0,00207. La deuxième est plus proche, ce qui illustre la stabilisation des fréquences. Aucune amélioration à chaque étape n’est garantie : un prochain succès ou échec peut augmenter l’écart.
\end{corrige}\end{blocLivre}
